\clearpage
\section{Worked outsider-facing translation of a BKM-type blow-up heuristic}

This appendix is optional and non-load-bearing. It is included only to make visible, in ordinary PDE
notation, one explicit instance of the mechanism used abstractly in the main proof: a familiar blow-up
heuristic is translated into a forbidden same-scope carrier move. Nothing in this appendix is needed
for the proof of Theorem~\ref{thm:official-periodic-pde-theorem}; its role is explanatory and audit-facing.
Appendix~D records the broader obstruction ledger; the present appendix is one concrete worked row of that ledger.

\subsection{A BKM-type blow-up observable}
Let $x\in I_{\mathrm{car}}(\T^3,\nu)$ and let
\[
\mathcal R_{\mathrm{NS}}(x)=(u_x,p_x)
\]
be its maximal realized trajectory on $[0,T_*(x))$. Define the BKM-type vorticity-growth observable
\[
\mathcal I_x(t):=\int_0^t \|\omega_x(s)\|_{L^\infty(\T^3)}\,\dd s,
\qquad 0\le t<T_*(x).
\]

\begin{remark}[Why this is the right test case]
The Beale--Kato--Majda criterion is a standard continuation criterion for smooth solutions: if
$\int_0^{T}\|\omega(s)\|_{L^\infty}\,\dd s<\infty$, then the solution continues past $T$.
Consequently, any finite-time breakdown hypothesis for a smooth trajectory forces the growth of the
observable $\mathcal I_x(t)$ as $t\uparrow T_*(x)$. This makes $\mathcal I_x$ the cleanest outsider-facing
example of a PDE-side quantity that appears to detect breakdown.
\end{remark}

\begin{lemma}[BKM growth at a finite breakdown horizon]\label{lem:bkm-growth-at-breakdown}
Let $(u_x,p_x)=\mathcal R_{\mathrm{NS}}(x)$ be a smooth realized trajectory on $[0,T_*(x))$ with
$T_*(x)<\infty$. Then
\[
\lim_{t\uparrow T_*(x)} \mathcal I_x(t)=\infty.
\]
\end{lemma}

\begin{proof}
If $\sup_{0\le t<T_*(x)}\mathcal I_x(t)<\infty$, then the Beale--Kato--Majda criterion applies on
$[0,T_*(x))$ and yields smooth continuation past $T_*(x)$, contradicting maximality of the realized
trajectory. Therefore $\mathcal I_x(t)$ cannot remain bounded as $t\uparrow T_*(x)$.
\end{proof}

\subsection{Threshold selectors and forbidden same-scope gates}
The structural point is not that $\mathcal I_x$ fails to exist. It is that using it as an independent
same-scope continuation test requires an extra threshold, tolerance, or horizon selector not contained
in the fixed carrier.

\begin{definition}[Finite-threshold BKM gate]\label{def:finite-threshold-bkm-gate}
Fix $M>0$ and $T>0$. The associated finite-threshold BKM gate on realized trajectories is
\[
G^{\mathrm{BKM}}_{M,T}(x):=\mathbf 1\!\left[\sup_{0\le t<T}\mathcal I_x(t)\le M\right].
\]
More generally, a BKM-profile gate is any same-scope rule asserting continuation from a comparison of
$\mathcal I_x(t)$ with a prescribed bound, threshold, tolerance, or growth profile as $t\uparrow T$.
\end{definition}

\begin{proposition}[Worked translation: BKM-type continuation tests force $S3$]\label{prop:bkm-forces-s3}
Let $x\in I_{\mathrm{car}}(\T^3,\nu)$ and suppose a same-scope continuation statement for the realized
trajectory $\mathcal R_{\mathrm{NS}}(x)$ is certified by a finite-threshold BKM gate or, more generally, by a
BKM-profile gate in the sense of Definition~\ref{def:finite-threshold-bkm-gate}. Then the certification
instantiates the forbidden move $S3$: it introduces a threshold, tolerance, selector, or horizon rule
not already licensed by the fixed carrier.
\end{proposition}

\begin{proof}
The observable $\mathcal I_x(t)$ is built from the realized vorticity field and is therefore a lawful
PDE-side quantity. The extra structure enters when one attempts to turn $\mathcal I_x$ into a
same-scope continuation criterion.

For the finite-threshold gate $G^{\mathrm{BKM}}_{M,T}$, the rule depends on the choice of the external
parameter $M$ and on the horizon parameter $T$. Neither parameter is determined by the carrier seed,
by the realization map, or by the horizon continuation discriminator already present in Part~III. The
same is true for any profile gate depending on a prescribed growth profile $\Phi(T-t)$, tolerance
$\varepsilon$, or selector determining what growth rate counts as admissible. Such data are not forced by
lawful redescription of the realized trajectory; they are additional same-scope choices.

Therefore the gate does not merely read continuation content already present on the carrier. It adds a
threshold / tolerance / selector rule for deciding continuation from the BKM observable. That is exactly
the structural role labeled $S3$ in the main text.
\end{proof}

\begin{corollary}[Concrete PDE-to-carrier translation of one blow-up heuristic]\label{cor:concrete-bkm-translation}
A BKM-style blow-up heuristic does not produce an admissible independent continuation mechanism on the
fixed torus carrier. Interpreted as a same-scope continuation gate, it forces the forbidden branch $S3$.
If one weakens the gate by allowing local exceptions or selective tolerance to large-vorticity episodes,
then one additionally enters the bundle-exemption branches $S1$ or $S2$.
\end{corollary}

\begin{proof}
The first statement is Proposition~\ref{prop:bkm-forces-s3}. If the gate is weakened by saying, for
example, that only selected portions of the trajectory, selected regions of the flow, or selected
instances of large vorticity matter, then the rule no longer treats the bundle clauses uniformly on the
fixed carrier. That is precisely the move-class captured by $S1$ or $S2$ in the main theorem package.
\end{proof}

\begin{remark}[What this appendix shows]
The main proof does not require any special analytic calculation attached to the BKM observable. The
role of this appendix is different: it shows, in one concrete and familiar PDE example, how a standard
blow-up heuristic is translated by hand into the carrier language. The contradiction in the main line is
therefore not a bare reclassification; it is the same mechanism exhibited abstractly in the proof spine,
made explicit here for one canonical continuation test.
\end{remark}



\clearpage
